\chapter{Debut Theorem}
\label{chap:debut_theorem}


\section{Choquet's capacitability theorem}

This section is devoted to the proof of Choquet's capacitability theorem, which is a key ingredient in the proof of the debut theorem. The presentation follows \cite{bichteler2002stochastic} for the definition of analytic sets (that definition changes a lot between sources) and for the general proof steps, and \cite{he2019semimartingale} and the PlanetMath website for many of the proofs of individual lemmas.
Although the definition of analytic sets in \cite{he2019semimartingale} is different, the proofs follow the same general steps and the same ideas, and that book is much more detailed than \cite{bichteler2002stochastic}.

\subsection{Pavings and Compact systems}

A paving is simply a set of sets.

\begin{definition}\label{def:image2_prod}
  \mathlibok
  % no \lean{}: it's Set.image2 of a product
The product of two pavings $S$ and $T$ is the set of sets of the form $s \times t$ where $s \in S$ and $t \in T$. We denote that product by $S \times T$.
\end{definition}


\begin{definition}\label{def:supClosure}
  \mathlibok
  \lean{supClosure}
For a paving $S$, we denote by $S^{\cup f}$ the set of finite unions of sets in $S$.
\end{definition}


\begin{definition}\label{def:infClosure}
  \mathlibok
  \lean{infClosure}
For a paving $S$, we denote by $S^{\cap f}$ the set of finite intersections of sets in $S$.
\end{definition}


\begin{definition}\label{def:countableSupClosure}
  \leanok
  \lean{countableSupClosure}
We denote by $S_\sigma$ the set of countable unions of sets in $S$.
\end{definition}


\begin{definition}\label{def:countableInfClosure}
  \leanok
  \lean{countableInfClosure}
We denote by $S_\delta$ the set of countable intersections of sets in $S$.
\end{definition}


\begin{lemma}\label{lem:InfClosed.mem_countableInfClosure_iff}
  \uses{def:countableInfClosure}
  \leanok
  \lean{InfClosed.mem_countableInfClosure_iff}
If a paving $S$ is closed under pairwise intersection, then a set $s$ is in $S_\delta$ if and only if there exists a monotone decreasing family of sets $(A_n)_{n\in\mathbb{N}}$ in $S$ such that $s = \bigcap_{n\in\mathbb{N}} A_n$.
\end{lemma}

\begin{proof}\leanok
The ``monotone decreasing'' condition is the only non-trivial part of the statement: without it it would just be the definition of $S_\delta$.
If $s \in S_\delta$ there exists a family $(B_n)$ that satisfies the condition but may not be monotone. We can then define $A_n = \bigcap_{k \le n} B_k$, and the family $(A_n)$ is monotone decreasing and satisfies $\bigcap_{n\in\mathbb{N}} A_n = \bigcap_{n\in\mathbb{N}} B_n = s$.
\end{proof}


\begin{lemma}\label{lem:SupClosed.mem_countableSupClosure_iff}
  \uses{def:countableSupClosure}
  \leanok
  \lean{SupClosed.mem_countableSupClosure_iff}
If a paving $S$ is closed under pairwise union, then a set $s$ is in $S_\sigma$ if and only if there exists a monotone increasing family of sets $(A_n)_{n\in\mathbb{N}}$ in $S$ such that $s = \bigcup_{n\in\mathbb{N}} A_n$.
\end{lemma}

\begin{proof}\leanok
Similar to the proof of Lemma~\ref{lem:InfClosed.mem_countableInfClosure_iff}.
\end{proof}


\begin{definition}\label{def:IsCompactSystem}
  \mathlibok
  \lean{IsCompactSystem}
A set of sets $K$ is a compact system if for every countable family $(C_n)_{n \in \mathbb{N}}$ of sets in $K$ such that $\bigcap_{n \in\mathbb{N}} C_n = \emptyset$, there exists a finite subset $S$ of $\mathbb{N}$ such that $\bigcap_{n \in S} C_n = \emptyset$.
\end{definition}


\begin{lemma}\label{lem:isCompactSystem_isCompact_isClosed}
  \uses{def:IsCompactSystem}
  \leanok
  \lean{isCompactSystem_isCompact_isClosed}
The family of compact closed sets of a topological space is a compact system.
\end{lemma}

\begin{proof}\leanok
  \uses{def:IsCompactSystem}

\end{proof}


\begin{lemma}\label{lem:isCompactSystem_isCompact}
  \uses{def:IsCompactSystem}
  \leanok
  \lean{isCompactSystem_isCompact}
The family of compact sets of a T2 topological space is a compact system.
\end{lemma}

\begin{proof}\leanok
  \uses{lem:isCompactSystem_isCompact_isClosed}
The T2 assumption ensures that compact sets are closed, so the family of compact sets is the family of compact closed sets, which is a compact system by Lemma~\ref{lem:isCompactSystem_isCompact_isClosed}.
\end{proof}


\begin{lemma}\label{lem:IsCompactSystem.mono}
  \uses{def:IsCompactSystem}
  \leanok
  \lean{IsCompactSystem.mono}
If $S \subseteq T$ and $T$ is a compact system, then $S$ is a compact system.
\end{lemma}

\begin{proof}\leanok
  \uses{def:IsCompactSystem}
Any sequence of sets in $S$ is a sequence of sets in $T$, so if the intersection of the sequence is empty, then there is a finite subset of the sequence whose intersection is empty by the compact system property of $T$.
\end{proof}


\begin{lemma}\label{lem:isCompactSystem_Icc}
  \uses{def:IsCompactSystem}
  \leanok
  \lean{isCompactSystem_Icc}
The family of closed compact intervals of $\mathbb{R}$ is a compact system.
\end{lemma}

\begin{proof}\leanok
  \uses{lem:isCompactSystem_isCompact, lem:IsCompactSystem.mono}
The family of closed compact intervals of $\mathbb{R}$ is a family of compact sets, so it is a compact system by Lemma~\ref{lem:isCompactSystem_isCompact} and~\ref{lem:IsCompactSystem.mono}.
\end{proof}


\begin{lemma}\label{lem:IsCompactSystem.pi}
  \uses{def:IsCompactSystem}
The product of a countable family of compact systems is a compact system.
\end{lemma}

\begin{proof}

\end{proof}


\begin{lemma}\label{lem:IsCompactSystem.sigma}
  \uses{def:IsCompactSystem}
The sum of a countable family of compact systems is a compact system.
\end{lemma}

\begin{proof}

\end{proof}


\begin{lemma}\label{lem:IsCompactSystem.image2_prod}
  \uses{def:IsCompactSystem}
  \leanok
  \lean{IsCompactSystem.image2_prod}
The product of two compact systems $S \times T$ is a compact system.
\end{lemma}

\begin{proof}

\end{proof}


\begin{lemma}\label{lem:IsCompactSystem.supClosure}
  \uses{def:IsCompactSystem, def:supClosure}
  \leanok
  \lean{IsCompactSystem.supClosure}
If $S$ is a compact system, then $S^{\cup f}$ is a compact system.
\end{lemma}

\begin{proof}

\end{proof}


\begin{lemma}\label{lem:IsCompactSystem.countableInfClosure}
  \uses{def:IsCompactSystem, def:countableInfClosure}
If $S$ is a compact system, then $S_\delta$ is a compact system.
\end{lemma}

\begin{proof}

\end{proof}


\begin{lemma}\label{lem:IsCompactSystem.infClosure}
  \uses{def:IsCompactSystem, def:infClosure}
  \leanok
  \lean{IsCompactSystem.infClosure}
If $S$ is a compact system, then $S^{\cap f}$ is a compact system.
\end{lemma}

\begin{proof}
  \uses{lem:IsCompactSystem.mono, lem:IsCompactSystem.countableInfClosure}
$S_\delta$ is a compact system by Lemma~\ref{lem:IsCompactSystem.countableInfClosure}, and $S^{\cap f} \subseteq S_\delta$, so we can apply Lemma~\ref{lem:IsCompactSystem.mono}.
\end{proof}


\begin{theorem}\label{thm:fst_iInter_of_supClosure_image2_prod_of_antitone}
  \uses{def:IsCompactSystem, def:supClosure, def:image2_prod}
  \leanok
  \lean{MeasureTheory.fst_iInter_of_supClosure_image2_prod_of_antitone}
Let $S$ be a paving of $\mathcal{X}$ and $T$ be a compact system of $\mathcal{K}$.
Let $(B_n)_{n \in \mathbb{N}}$ be a monotone decreasing sequence of sets in $(S \times T)^{\cup f}$.
Let $\pi_{\mathcal{X}} : \mathcal{X} \times \mathcal{K} \to \mathcal{X}$ be the projection on $\mathcal{X}$.
Then
\begin{align*}
  \pi_{\mathcal{X}}\left(\bigcap_{n\in\mathbb{N}} B_n\right) = \bigcap_{n\in\mathbb{N}} \pi_{\mathcal{X}}(B_n).
\end{align*}
\end{theorem}

\begin{proof}\leanok
  \uses{lem:IsCompactSystem.supClosure}

\end{proof}




\subsection{Analytic sets}

\begin{definition}[Analytic set]\label{def:IsPavingAnalytic}
  \uses{def:IsCompactSystem, def:countableSupClosure, def:countableInfClosure, def:image2_prod}
  \leanok
  \lean{MeasureTheory.IsPavingAnalytic, MeasureTheory.IsPavingAnalyticFor}
Let $F$ be a paving of a type $\mathcal{X}$.
A set $s$ of $\mathcal{X}$ is said to be $F$-analytic if there exists a nonempty type $\mathcal{K}$ and a compact system $K$ of $\mathcal{K}$ with $\emptyset \in K$ such that $s$ is the projection on $\mathcal{X}$ of a set in $(F \times K)_{\sigma \delta}$.

We denote the set of $F$-analytic sets by $\mathcal{A}(F)$.
\end{definition}

Lean note: we can't quantify over universes, so we have to restrict $\mathcal{K}$ to \texttt{Type}. We have an auxiliary definition \texttt{IsPavingAnalyticFor} that specifies which $\mathcal{K}$ is used.


\begin{lemma}\label{lem:isPavingAnalytic_of_prop}
  \uses{def:IsPavingAnalytic}
  \leanok
  \lean{MeasureTheory.isPavingAnalytic_of_mem}
Every set in $F$ is $F$-analytic: $F \subseteq \mathcal{A}(F)$.
\end{lemma}

\begin{proof}\leanok

\end{proof}


\begin{lemma}\label{lem:IsPavingAnalytic.mono}
  \uses{def:IsPavingAnalytic}
  \leanok
  \lean{MeasureTheory.IsPavingAnalytic.mono}
If $F \subseteq G$ and $s$ is $F$-analytic, then $s$ is $G$-analytic: $\mathcal{A}(F) \subseteq \mathcal{A}(G)$.
\end{lemma}

\begin{proof}\leanok

\end{proof}


\begin{lemma}\label{lem:IsPavingAnalytic.exists_mem_countableSupClosure_superset}
  \uses{def:IsPavingAnalytic}
  \leanok
  \lean{MeasureTheory.IsPavingAnalyticFor.exists_mem_countableSupClosure_superset}
If $s$ is $F$-analytic, then there exists a set $t$ in $F_\sigma$ such that $s \subseteq t$.
\end{lemma}

\begin{proof}\leanok

\end{proof}


\begin{lemma}\label{lem:IsPavingAnalytic.iUnion}
  \uses{def:IsPavingAnalytic}
  \leanok
  \lean{MeasureTheory.IsPavingAnalytic.iUnion, MeasureTheory.IsPavingAnalytic.union}
A countable union of $F$-analytic sets is $F$-analytic.
\end{lemma}

\begin{proof}\leanok
  \uses{lem:IsCompactSystem.pi}

\end{proof}


\begin{lemma}\label{lem:IsPavingAnalytic.iInter}
  \uses{def:IsPavingAnalytic}
  \leanok
  \lean{MeasureTheory.IsPavingAnalytic.iInter, MeasureTheory.IsPavingAnalytic.inter}
A countable intersection of $F$-analytic sets is $F$-analytic.
\end{lemma}

\begin{proof}\leanok
  \uses{lem:IsCompactSystem.sigma}

\end{proof}


\begin{lemma}\label{lem:IsPavingAnalytic.fst}
  \uses{def:IsPavingAnalytic, def:IsCompactSystem}
  \leanok
  \lean{MeasureTheory.IsPavingAnalytic.fst}
Let $F$ be a set of sets of $\mathcal{X}$ and $K$ be a compact system of $\mathcal{K}$, with $\emptyset \in K$.
If a set $s$ of $\mathcal{X} \times \mathcal{K}$ is $F \times K$-analytic, then its projection on $\mathcal{X}$ is $F$-analytic.
\end{lemma}

\begin{proof}\leanok
  \uses{lem:IsCompactSystem.image2_prod}

\end{proof}


\begin{lemma}\label{lem:IsPavingAnalytic.prod_left}
  \uses{def:IsPavingAnalytic}
  \leanok
  \lean{MeasureTheory.IsPavingAnalytic.prod_left}
If $t \in T$ and $s$ is $F$-analytic, then $t \times s$ is $(T \times F)$-analytic: $T \times \mathcal{A}(F) \subseteq \mathcal{A}(T \times F)$.
\end{lemma}

\begin{proof}\leanok

\end{proof}


\begin{lemma}\label{lem:IsPavingAnalytic.prod}
  \uses{def:IsPavingAnalytic}
  \leanok
  \lean{MeasureTheory.IsPavingAnalytic.prod}
If $s$ is $F$-analytic and $t$ is $G$-analytic, then $s \times t$ is $(F \times G)$-analytic: $\mathcal{A}(F) \times \mathcal{A}(G) \subseteq \mathcal{A}(F \times G)$.
\end{lemma}

\begin{proof}\leanok
  \uses{lem:IsPavingAnalytic.prod_left, lem:IsPavingAnalytic.exists_mem_countableSupClosure_superset}

\end{proof}


\begin{lemma}\label{lem:isPavingAnalytic_isPavingAnalytic}
  \uses{def:IsPavingAnalytic}
  \leanok
  \lean{MeasureTheory.isPavingAnalytic_isPavingAnalytic}
$\mathcal{A}(\mathcal{A}(F)) = \mathcal{A}(F)$.
\end{lemma}

\begin{proof}\leanok
  \uses{lem:IsPavingAnalytic.iInter, lem:IsPavingAnalytic.iUnion, lem:IsPavingAnalytic.prod, lem:IsPavingAnalytic.fst}

\end{proof}



\subsection{Analytic sets in measurable spaces}

We denote by $K_{\mathbb{R}}$ the paving of compact sets of $\mathbb{R}$ and let $I_{\mathbb{R}}$ be the paving of closed compact intervals of $\mathbb{R}$.
For a measurable space $\mathcal{X}$, we denote by $M_{\mathcal{X}}$ the paving of measurable sets of $\mathcal{X}$.

\begin{lemma}\label{lem:isPavingAnalytic_of_measurableSet_generateFrom}
  \uses{def:IsPavingAnalytic}
  \leanok
  \lean{MeasureTheory.isPavingAnalytic_of_measurableSet_generateFrom}
Let $F$ be a paving with $\emptyset \in F$ and let $\sigma(F)$ be the $\sigma$-algebra generated by $F$.
Suppose that for every $t \in F$, $t^c$ is $F$-analytic.
Then every set in $\sigma(F)$ is $F$-analytic: $\sigma(F) \subseteq \mathcal{A}(F)$.
\end{lemma}

\begin{proof}\leanok
  \uses{lem:isPavingAnalytic_of_prop, lem:IsPavingAnalytic.iUnion, lem:IsPavingAnalytic.iInter}

\end{proof}


\begin{lemma}\label{lem:MeasurableSet.isPavingAnalytic_Icc_real}
  \uses{def:IsPavingAnalytic}
  \leanok
  \lean{MeasurableSet.isPavingAnalytic_Icc_real}
A measurable set of $\mathbb{R}$ is $I_{\mathbb{R}}$-analytic.
That is, $M_{\mathbb{R}} \subseteq \mathcal{A}(I_{\mathbb{R}})$.
\end{lemma}

\begin{proof}\leanok
  \uses{lem:isPavingAnalytic_of_measurableSet_generateFrom, lem:IsPavingAnalytic.iUnion, lem:isPavingAnalytic_of_prop}

\end{proof}


\begin{lemma}\label{lem:IsPavingAnalytic_measurableSet_iff_isPavingAnalytic_Icc}
  \uses{def:IsPavingAnalytic}
  \leanok
  \lean{MeasureTheory.IsPavingAnalytic_measurableSet_iff_isPavingAnalytic_Icc}
A set of $\mathbb{R}$ is $M_{\mathbb{R}}$-analytic if and only if it is $I_{\mathbb{R}}$-analytic.
\end{lemma}

\begin{proof}\leanok
  \uses{lem:isPavingAnalytic_isPavingAnalytic, lem:MeasurableSet.isPavingAnalytic_Icc_real}

\end{proof}


\begin{lemma}\label{lem:MeasurableSet.isPavingAnalytic_image2_prod}
  \uses{def:IsPavingAnalytic}
  \leanok
  \lean{MeasurableSet.isPavingAnalytic_image2_prod}
Let $\mathcal{X}$ be a measurable space and $s$ be a measurable set of $\mathcal{X} \times \mathbb{R}$.
Then $s$ is $(M_{\mathcal{X}} \times I_{\mathbb{R}})$-analytic.
\end{lemma}

\begin{proof}\leanok
  \uses{lem:isPavingAnalytic_of_measurableSet_generateFrom, lem:IsPavingAnalytic.iUnion, lem:isPavingAnalytic_of_prop}

\end{proof}


\begin{lemma}\label{lem:isPavingAnalytic_fst_of_image2_prod_measurableSet_Icc}
  \uses{def:IsPavingAnalytic}
  \leanok
  \lean{MeasureTheory.isPavingAnalytic_fst_of_image2_prod_measurableSet_Icc}
Let $\mathcal{X}$ be a measurable space and $s$ be a set of $\mathcal{X} \times \mathbb{R}$ which is $(M_{\mathcal{X}} \times I_{\mathbb{R}})$-analytic.
Then the projection $\pi_{\mathcal{X}}(s)$ is $M_{\mathcal{X}}$-analytic.
\end{lemma}

\begin{proof}\leanok
  \uses{lem:IsPavingAnalytic.fst, lem:isCompactSystem_Icc}

\end{proof}


\begin{lemma}\label{lem:MeasurableSet.isPavingAnalytic_fst}
  \uses{def:IsPavingAnalytic}
  \leanok
  \lean{MeasurableSet.isPavingAnalytic_fst}
Let $\mathcal{X}$ be a measurable space and $s$ be a measurable set of $\mathcal{X} \times \mathbb{R}$.
Then the projection of $s$ on $\mathcal{X}$ is $M_{\mathcal{X}}$-analytic.
\end{lemma}

\begin{proof}\leanok
  \uses{lem:MeasurableSet.isPavingAnalytic_image2_prod, lem:isPavingAnalytic_fst_of_image2_prod_measurableSet_Icc}

\end{proof}


\begin{definition}\label{def:IsMeasurableAnalytic}
  \leanok
  \lean{MeasureTheory.IsMeasurableAnalytic, MeasureTheory.IsMeasurableAnalyticFor}
We say that a set $s$ of a measurable space $\mathcal{X}$ is measurably analytic for a measurable space $\mathcal{K}$ if it is the projection of a measurable set of $\mathcal{X} \times \mathcal{K}$.

We say that $s$ is measurably analytic if it is measurably analytic for $\mathbb{R}$.
\end{definition}


\begin{lemma}\label{lem:IsMeasurableAnalyticFor.isMeasurableAnalytic}
  \uses{def:IsMeasurableAnalytic}
  \leanok
  \lean{MeasureTheory.IsMeasurableAnalyticFor.isMeasurableAnalytic}
Let $\mathcal{K}$ be a standard Borel space.
If a set $s$ of $\mathcal{X}$ is measurably analytic for $\mathcal{K}$, then it is measurably analytic.
\end{lemma}

\begin{proof}\leanok
There is a measurable embedding of $\mathcal{K}$ into $\mathbb{R}$.
\end{proof}


\begin{lemma}\label{lem:IsMeasurableAnalytic.isPavingAnalytic}
  \uses{def:IsMeasurableAnalytic, def:IsPavingAnalytic}
  \leanok
  \lean{MeasureTheory.IsMeasurableAnalytic.isPavingAnalytic}
If a set $s$ of $\mathcal{X}$ is measurably analytic, then it is $M_{\mathcal{X}}$-analytic.
\end{lemma}

\begin{proof}\leanok
  \uses{lem:MeasurableSet.isPavingAnalytic_fst}

\end{proof}



\subsection{Capacities and capacitable sets}

\begin{definition}\label{def:Capacity}
  \leanok
  \lean{MeasureTheory.Capacity}
Let $F$ be a set of sets of a type $\mathcal{X}$.
An $F$-capacity is a function $I$ from the sets of $\mathcal{X}$ to $\mathbb{R}_{+,\infty}$ such that
\begin{enumerate}
  \item $I$ is monotone for set inclusion: for $s \subseteq t$, $I(s) \le I(t)$,
  \item if $(s_n)_{n\in\mathbb{N}}$ is an increasing sequence of sets of $\mathcal{X}$, then $I(s_n)$ tends to $I(\bigcup_{n\in\mathbb{N}} s_n)$ at infinity,
  \item if $(s_n)_{n\in\mathbb{N}}$ is a decreasing sequence of sets in $F$, then $I(s_n)$ tends to $I(\bigcap_{n\in\mathbb{N}} s_n)$ at infinity.
\end{enumerate}
\end{definition}


\begin{lemma}\label{lem:Capacity.comp_fst_aux}
  \uses{def:Capacity, def:IsCompactSystem}
  \leanok
  \lean{MeasureTheory.Capacity.comp_fst}
Let $F$ be a set of sets of $\mathcal{X}$ and let $I$ be an $F$-capacity.
Suppose that $\emptyset \in F$ and $F$ is closed under pairwise unions.
Let $K$ be a compact system of $\mathcal{K}$.
Let $\pi_{\mathcal{X}} : \mathcal{X} \times \mathcal{K} \to \mathcal{X}$ be the projection on the first coordinate.
Then the function $I_{\mathcal{X}}(u) = I(\pi_{\mathcal{X}}(u))$ is a capacity for $(F \times K)^{\cup f}$.
\end{lemma}

\begin{proof}\leanok
  \uses{thm:fst_iInter_of_supClosure_image2_prod_of_antitone}

\end{proof}


\begin{definition}\label{def:Capacity.comp_fst}
  \uses{def:Capacity, def:IsCompactSystem}
  \leanok
  \lean{MeasureTheory.Capacity.comp_fst}
Let $F$ be a set of sets of $\mathcal{X}$ and let $I$ be an $F$-capacity.
Suppose that $\emptyset \in F$ and $F$ is closed under pairwise unions.
Let $K$ be a compact system of $\mathcal{K}$.
Let $\pi_{\mathcal{X}} : \mathcal{X} \times \mathcal{K} \to \mathcal{X}$ be the projection on the first coordinate.
Then we define a capacity for $(F \times K)^{\cup f}$ by $I_{\mathcal{X}}(u) = I(\pi_{\mathcal{X}}(u))$ (it is indeed a capacity by Lemma~\ref{lem:Capacity.comp_fst_aux}).
\end{definition}


\begin{definition}\label{def:IsCapacitable}
  \uses{def:Capacity}
  \leanok
  \lean{MeasureTheory.IsCapacitable}
Let $F$ be a set of sets of $\mathcal{X}$ and let $I$ be an $F$-capacity.
A set $s$ of $\mathcal{X}$ is said to be $I$-capacitable if for all $a < I(s)$, there exists a set $t \in F_\delta$ such that $t \subseteq s$ and $a \le I(t)$.
\end{definition}


\begin{lemma}\label{lem:isCapacitable_of_prop}
  \uses{def:IsCapacitable}
  \leanok
  \lean{MeasureTheory.isCapacitable_of_mem}
For $I$ an $F$-capacity, every set in $F$ is $I$-capacitable.
\end{lemma}

\begin{proof}\leanok

\end{proof}


\begin{lemma}\label{lem:isCapacitable_mem_countableInfClosure_countableSupClosure}
  \uses{def:IsCapacitable, def:Capacity, def:countableInfClosure, def:countableSupClosure}
  \leanok
  \lean{MeasureTheory.isCapacitable_mem_countableInfClosure_countableSupClosure}
Suppose that $F$ contains the empty set and is closed under pairwise intersections and unions, and let $I$ be an $F$-capacity.
Then every set of $F_{\sigma \delta}$ is $I$-capacitable.
\end{lemma}

\begin{proof}
  \uses{lem:SupClosed.mem_countableSupClosure_iff}

\end{proof}


\begin{lemma}\label{lem:mem_countableInfClosure_fst}
  \uses{def:countableInfClosure, def:image2_prod, def:IsCompactSystem}
  \leanok
  \lean{MeasureTheory.mem_countableInfClosure_fst}
Suppose that $F$ contains the empty set and is closed under pairwise intersections and unions, and let $I$ be an $F$-capacity.
Let $K$ be a compact system which is also closed under pairwise intersections.
Let $s \in ((F \times K)^{\cup f})_\delta$.
Then $\pi_{\mathcal{X}}(s) \in F_\delta$.
\end{lemma}

\begin{proof}\leanok
  \uses{thm:fst_iInter_of_supClosure_image2_prod_of_antitone, lem:InfClosed.mem_countableInfClosure_iff}

\end{proof}


\begin{lemma}\label{lem:IsCapacitable.fst}
  \uses{def:IsCapacitable, def:Capacity, def:Capacity.comp_fst, def:IsCompactSystem}
  \leanok
  \lean{MeasureTheory.IsCapacitable.fst}
Let $F$ be a set of sets of $\mathcal{X}$ and let $I$ be an $F$-capacity.
Suppose that $F$ contains the empty set and is closed under pairwise intersections and unions.
Let $K$ be a compact system of $\mathcal{K}$ which is also closed under pairwise intersections.
If a set $s$ of $\mathcal{X} \times \mathcal{K}$ is capacitable with respect to the capacity $I_{\mathcal{X}}$ of Definition~\ref{def:Capacity.comp_fst}, then its projection on $\mathcal{X}$ is capacitable with respect to $I$.
\end{lemma}

\begin{proof}\leanok
  \uses{lem:mem_countableInfClosure_fst}

\end{proof}


\begin{theorem}[Choquet's capacitability theorem]\label{thm:IsPavingAnalytic.isCapacitable}
  \uses{def:IsPavingAnalytic, def:IsCapacitable}
  \leanok
  \lean{MeasureTheory.IsPavingAnalytic.isCapacitable}
Let $F$ be a set of sets of $\mathcal{X}$ and let $I$ be an $F$-capacity.
Suppose that $F$ contains the empty set and is closed under pairwise intersections and unions.
Then every $F$-analytic set is $I$-capacitable.
\end{theorem}

\begin{proof}\leanok
  \uses{lem:isCapacitable_mem_countableInfClosure_countableSupClosure, lem:IsCompactSystem.infClosure, lem:IsCapacitable.fst}

\end{proof}



\subsection{Capacities and measures}


\begin{definition}\label{def:Measure.capacity}
  \uses{def:Capacity}
  \leanok
  \lean{MeasureTheory.Measure.capacity}
Every finite measure defines a capacity by $I(s) = \mu(s)$ (where $\mu(s)$ is the value of the outer measure defined by $\mu$, but in Lean that's how $\mu(s)$ is defined already).
\end{definition}


\begin{lemma}\label{lem:isCapacitable_measure_iff}
  \uses{def:IsCapacitable, def:Measure.capacity}
  \leanok
  \lean{MeasureTheory.isCapacitable_measure_iff}
A set is capacitable with respect to the capacity defined by a finite measure if and only if it is null-measurable for this measure.
\end{lemma}

\begin{proof}\leanok

\end{proof}


\begin{lemma}\label{lem:IsPavingAnalytic.nullMeasurableSet}
  \uses{def:IsPavingAnalytic}
  \leanok
  \lean{MeasureTheory.IsPavingAnalytic.nullMeasurableSet}
An $M_{\mathcal{X}}$-analytic set of a measurable space $\mathcal{X}$ is universally measurable: it is null-measurable for every finite measure on $\mathcal{X}$.
\end{lemma}

\begin{proof}\leanok
  \uses{thm:IsPavingAnalytic.isCapacitable, def:Measure.capacity, lem:isCapacitable_measure_iff}

\end{proof}


\begin{lemma}\label{lem:IsMeasurableAnalytic.nullMeasurableSet}
  \uses{def:IsMeasurableAnalytic}
  \leanok
  \lean{MeasureTheory.IsMeasurableAnalytic.nullMeasurableSet}
A measurably analytic set of a measurable space $\mathcal{X}$ is universally measurable: it is null-measurable for every finite measure on $\mathcal{X}$.
\end{lemma}

\begin{proof}\leanok
  \uses{lem:IsMeasurableAnalytic.isPavingAnalytic, lem:IsPavingAnalytic.nullMeasurableSet}

\end{proof}


\begin{theorem}[Measurable projection]\label{thm:MeasurableSet.nullMeasurableSet_fst}
  \uses{def:IsMeasurableAnalytic}
  \leanok
  \lean{MeasurableSet.nullMeasurableSet_fst}
Let $\mathcal{X}$ and $\mathcal{Y}$ be measurable spaces, with $\mathcal{Y}$ standard Borel.
Then the projection of a measurable set of $\mathcal{X} \times \mathcal{Y}$ on $\mathcal{X}$ is universally measurable.
\end{theorem}

\begin{proof}\leanok
  \uses{lem:IsMeasurableAnalyticFor.isMeasurableAnalytic, lem:IsMeasurableAnalytic.nullMeasurableSet}

\end{proof}




\section{Monotone class theorem}


\begin{definition}[Monotone class]\label{def:monotone_class}
  \leanok
  \lean{MeasureTheory.MonotoneClass}
Let $\mathcal{M}$ be a collection of subsets of a set $X$. We say that $\mathcal{M}$ is a monotone class if it is closed under countable monotone unions and countable monotone intersections, i.e.:
\begin{enumerate}
  \item if \( A_1, A_2, \ldots \in M \) and \( A_1 \subseteq A_2 \subseteq \cdots \), then
  \( \bigcup_{i=1}^\infty A_i \in M \),
  \item if \( B_1, B_2, \ldots \in M \) and \( B_1 \supseteq B_2 \supseteq \cdots \), then
  \( \bigcap_{i=1}^\infty B_i \in M \).
\end{enumerate}
Given a collection $\mathcal{F}$ of subsets of $X$, we call the smallest monotone class containing $\mathcal{F}$ the monotone class generated by $\mathcal{F}$.
\end{definition}


\begin{theorem}[Monotone class theorem]\label{thm:monotone_class}
  \uses{def:monotone_class}
  \leanok
  \lean{MeasureTheory.MonotoneClass.generateFrom_eq}
Let \(G\) be an algebra of subsets of a set \(X\). Then the monotone class generated by \(G\) coincides with the $\sigma$-algebra generated by \(G\).
\end{theorem}

\begin{proof}\leanok

\end{proof}




\section{Debut Theorem}


\begin{definition}[Progressively measurable set]\label{def:progr_meas_set}
  \leanok
  \uses{def:ProgMeasurable}
  \lean{MeasureTheory.ProgMeasurableSet}
A subset of $[0, \infty) \times \Omega$ is progressively measurable if its indicator is a progressively measurable process.
\end{definition}


\begin{definition}[Debut of a set]\label{def:debut_set}
  \leanok
  \lean{MeasureTheory.debut}
Let $E \subseteq{} [0, \infty) \times \Omega $, define $D_E = \inf\left\lbrace t \geq 0\ :\ (t, \omega) \in E\right\rbrace$, the debut of $E$.
\end{definition}


\begin{theorem}\label{thm:debut_of_progr_meas_is_stop_time}
  \leanok
  \lean{MeasureTheory.isStoppingTime_debut}
  \uses{def:progr_meas_set, def:debut_set, def:IsStoppingTime, def:hasUsualConditions}
If $E$ is a progressively measurable set and the filtration satisfies the usual conditions, then $D_E$ is a stopping time.
\end{theorem}

\begin{proof}\leanok
  \uses{thm:MeasurableSet.nullMeasurableSet_fst}

\end{proof}




\section{Hitting times}

\begin{theorem}\label{thm:hitting_is_stopping_time}
  \uses{def:IsStoppingTime, def:hittingAfter, def:ProgMeasurable, def:hasUsualConditions}
  \leanok
  \lean{MeasureTheory.isStoppingTime_hittingAfter'}
If $X : T \to \Omega \to E$ is a progressively measurable process with respect to a filtration satisfying the usual conditions and $B$ is a Borel-measurable subset of $E$, then the hitting time of $X$ in $B$ is a stopping time.
\end{theorem}

\begin{proof}\leanok
  \uses{thm:debut_of_progr_meas_is_stop_time}
Since $B$ is a Borel subset of $\mathcal{S}$ and $X$ is progressively measurable, then $\mathbf{1}_B(X_t)$ is also progressively measurable.
The hitting time is then the debut of the set $E = \{(s,\omega) : \mathbf{1}_B(X_s(\omega)) = 1\}$, and therefore is a stopping time by Theorem~\ref{thm:debut_of_progr_meas_is_stop_time}.
\end{proof}


\begin{definition}\label{def:leastGE}
  \mathlibok
  \lean{MeasureTheory.leastGE}
Let $T$ be a conditionally complete linear order with a bottom element and let $R$ be a preorder.
For a process $X : T \to Ω \to R$ and $a \in R$, define the random time
\begin{align*}
  \tau_{X \ge a} = \inf\{t \in T \mid X_t \ge a\} \: ,
\end{align*}
in which the infimum is infinite if the set is empty.
\end{definition}


\begin{lemma}\label{lem:isStoppingTime_leastGE}
  \uses{def:IsStoppingTime, def:leastGE}
  \leanok
  \lean{MeasureTheory.isStoppingTime_leastGE}
Let $T$ be a conditionally complete linear order with a bottom element, which is a Polish space for its order topology, on which we put the Borel $\sigma$-algebra.
Let $R$ be a preorder with a metrizable topology for which sets of the form $\{x \in R \mid x \ge a\}$ are closed, with its Borel $\sigma$-algebra.
If $X : T \to Ω \to R$ is a progressively measurable process with respect to a filtration satisfying the usual conditions, then for any $a \in R$, the random time $\tau_{X \ge a}$ is a stopping time.
\end{lemma}

\begin{proof}\leanok
  \uses{thm:hitting_is_stopping_time}
This is a direct application of Theorem~\ref{thm:hitting_is_stopping_time} with the set $B = [a, +\infty)$.
\end{proof}


\begin{corollary}\label{cor:isStoppingTime_leastGE_of_rightContinuous}
  \uses{def:IsStoppingTime, def:leastGE, def:rightContinuous, def:adapted, def:RightContinuous}
If $X : ι \to Ω \to ℝ$ is a right-continuous and adapted process with respect to a filtration satisfying the usual conditions, then for any $a \in \mathbb{R}$, the random time $\tau_{X \ge a}$ is a stopping time.
\end{corollary}

\begin{proof}
  \uses{lem:isStoppingTime_leastGE, lem:Adapted.progMeasurable_of_rightContinuous}
This follows from Lemma~\ref{lem:isStoppingTime_leastGE} since $X$ is progressively measurable by Lemma~\ref{lem:Adapted.progMeasurable_of_rightContinuous}.
\end{proof}


\begin{definition}\label{def:leastGT}
  \leanok
  \lean{MeasureTheory.leastGT}
Let $T$ be a conditionally complete linear order with a bottom element and let $R$ be a preorder.
For a process $X : T \to Ω \to R$ and $a \in R$, define the random time
\begin{align*}
  \tau_{X > a} = \inf\{t \in T \mid X_t > a\} \: ,
\end{align*}
in which the infimum is infinite if the set is empty.
\end{definition}


\begin{lemma}\label{lem:isStoppingTime_leastGT}
  \uses{def:IsStoppingTime, def:leastGT}
  \leanok
  \lean{MeasureTheory.isStoppingTime_leastGT}
Let $T$ be a conditionally complete linear order with a bottom element, which is a Polish space for its order topology, on which we put the Borel $\sigma$-algebra.
Let $R$ be a preorder with a metrizable topology for which sets of the form $\{x \in R \mid x \ge a\}$ are closed, with its Borel $\sigma$-algebra.
If $X : T \to Ω \to R$ is a progressively measurable process with respect to a filtration satisfying the usual conditions, then for any $a \in R$, the random time $\tau_{X > a}$ is a stopping time.
\end{lemma}

\begin{proof}\leanok
  \uses{thm:hitting_is_stopping_time}
This is a direct application of Theorem~\ref{thm:hitting_is_stopping_time} with the set $B = (a, +\infty)$.
\end{proof}


\begin{lemma}\label{lem:leastGT_lt_iff}
  \uses{def:leastGT}
  \leanok
  \lean{MeasureTheory.leastGT_lt_iff}
For $t \in T$, $\tau_{X > a} < t$ if and only if there exists $s < t$ such that $X_s > a$.
\end{lemma}

\begin{proof}\leanok
  \uses{def:leastGT}

\end{proof}
